\section{Problema educacional e justificativa}

Na rotina escolar, alguns estudantes recorrem à inteligência artificial para obter resultados imediatos, sem necessariamente se apropriar do conhecimento envolvido. Quando esse uso se resume a buscar respostas prontas, o funcionamento da ferramenta fica oculto para quem a utiliza. Nesse cenário, ensinar conceitos básicos de \textit{machine learning} ajuda o jovem a conhecer o processo por trás das respostas que recebe. Soma-se a isso a falta de oportunidades de trabalho com dados. Em geral, tabelas e gráficos apresentados no Ensino Médio chegam prontos, inseridos em questões. Nessas situações, a leitura dos valores pode ocorrer sem uma discussão sobre a origem dos dados ou sobre o que mudaria no resultado se um valor ou uma condição fosse diferente.

Compreender um modelo de \textit{machine learning} exige algum contato prévio com dados, além de noções básicas de estatística e probabilidade. Já existem trabalhos voltados ao ensino do tema na Educação Básica. Em um mapeamento de 30 unidades didáticas, \citeonline{marques2020teaching} observaram que várias delas cobrem apenas os processos mais acessíveis, como a gestão de dados, ou apresentam o aprendizado e o teste dos modelos de forma abstrata, deixando parte do funcionamento interno oculto. Isso sugere que, nessas unidades, o contato com dados pode se restringir à gestão, sem mostrar como os dados e os parâmetros alteram o resultado do modelo. No Ensino Médio, a revisão de \citeonline{martins2023findings} indica que os estudantes conseguem compreender conceitos básicos de \textit{machine learning} e aplicá-los em tarefas com algoritmos. A mesma revisão aponta a necessidade de avaliações mais rigorosas sobre esse ensino. A presença da inteligência artificial na escola também pede atividades que vinculem os dados e a finalidade da tarefa ao modo como regras e parâmetros alteram o resultado do modelo.

Quando os dados chegam prontos e as relações entre eles não são trabalhadas, o estudante tende a aprender os conceitos de \textit{machine learning} apenas pelo nome. O vocabulário básico da área, que vai do treinamento do modelo ao ajuste de parâmetros para gerar previsões, acaba repetido sem conexão real. Um estudante pode dizer que um modelo foi treinado sem saber explicar o que muda quando um parâmetro é alterado, e pode reconhecer uma reta no gráfico sem relacioná-la a uma previsão. O problema educacional está na passagem entre nomear um conceito e utilizá-lo para explicar uma situação.

Há ainda uma diferença importante entre apresentar a inteligência artificial como tema de cultura digital e ensiná-la como conteúdo de Computação. A discussão cultural é indispensável para problematizar impactos sociais, abrangendo desde vieses e riscos até questões de autoria. No ensino de Computação, o estudante também precisa investigar, na prática, uma representação de dados e uma regra computacional. A proposta não reduz questões éticas a uma equação ou a um gráfico. Considera-se que a discussão sobre limites passa a ter outra base quando o estudante reconhece que modelos resultam de decisões humanas de modelagem.

Romper com a recepção passiva de resultados prontos exige que o estudante faça previsões antes de alterar dados ou parâmetros para observar e explicar o efeito produzido. Esse processo começa antes do modelo, na origem dos dados. Um estudante que monta uma tabela com medidas da própria turma sabe de onde vêm os números e o que cada coluna representa. Essa experiência torna a leitura do gráfico menos mecânica. Em seguida, ver uma reta se deslocar quando os dados mudam, ou a mudança de categoria quando o limiar é alterado, produz evidências visíveis. Ao explicar o que observou, o estudante pode relacionar o efeito visual à alteração que realizou. Essa organização se aproxima das abordagens práticas, baseadas em problemas ou projetos, que \citeonline{martins2023findings} identificaram como bem-sucedidas no Ensino Médio. Compreender de onde vêm os dados e o que muda com um parâmetro é o que permite ao estudante questionar o resultado que recebe, e não apenas aceitá-lo.

O produto educacional prevê uma plataforma interativa em português brasileiro, articulada a uma sequência didática para estudantes do Ensino Médio. A plataforma apresenta cinco módulos introdutórios: regressão linear, regressão logística, \textit{k}-vizinhos mais próximos, árvore de decisão e \textit{k}-means. Os modelos serão tratados como conteúdos centrais, sem receber necessariamente a mesma formalização matemática. Com a regressão linear, a turma analisa a linha de tendência para fazer previsões e medir o erro, passando depois à probabilidade e ao limiar na regressão logística. Já o \textit{k}-NN destaca distância e votação, a árvore de decisão explicita regras e o \textit{k}-means agrupa dados não rotulados.

O uso da plataforma, por si só, não será tomado como evidência de aprendizagem. Ferramentas visuais podem apoiar a compreensão, mas seu potencial depende de como as representações são desenhadas e das tarefas cognitivas pedidas aos estudantes \cite{ainsworth2006deft}. Por isso, cada módulo partirá de uma situação-problema e da origem dos dados para que a turma levante uma hipótese antes da manipulação controlada. A atividade avança, assim, da previsão inicial ao experimento prático, culminando na explicação do que foi observado. A devolutiva do experimento participa da mediação docente, aproximando-se da avaliação formativa discutida por \citeonline{black1998assessment}. Serão oito encontros de cerca de 50 minutos. No segundo, a turma poderá organizar uma tabela simples antes de estudar regressão linear; nos demais módulos, trabalhará com cenários sintéticos fornecidos pelo site. Um estudante que modifica o limiar de uma classificação precisa discutir por que a mudança ocorreu. Sem essa discussão, a manipulação do controle pode ficar desvinculada do conceito que se pretende ensinar.

A justificativa profissional está na necessidade de um recurso aplicável por professores da Educação Básica em turmas heterogêneas, considerando a limitação de tempo e a variação de infraestrutura entre escolas. Em um programa profissional como o ProfComp, a relevância do estudo reside em articular um problema escolar ao conteúdo de Computação por meio da intervenção docente com um produto utilizável. A análise da plataforma considera seu uso em aula, pois o estudo descarta a ideia de que uma interface interativa isolada resolva dificuldades de aprendizagem.

Espera-se contribuir para aproximar a inteligência artificial do Ensino Médio regular, mantendo visível como cada modelo funciona sem exigir o aprofundamento próprio dos cursos técnicos. A contribuição é localizada e dependerá das condições de aplicação. Documentar essas condições pode ajudar outros docentes a avaliar se uma proposta semelhante serve ao seu contexto, sem antecipar generalizações que a pesquisa não poderá sustentar.

\section{Pergunta de pesquisa}

Que mudanças e permanências podem ser identificadas nas explicações de estudantes do Ensino Médio sobre fundamentos introdutórios de \textit{machine learning} ao longo de uma sequência didática apoiada por uma plataforma interativa?

A pergunta dirige a análise para o modo como os estudantes justificam escolhas e explicam gráficos ou tabelas em diferentes momentos da sequência. Em vez de medir o efeito isolado da plataforma, o estudo examinará o que muda ou permanece nas dificuldades de relacionar os dados e a tarefa ao ajuste de parâmetros na regra do modelo. As produções serão elaboradas ao longo de um percurso que combina o uso individual da plataforma à mediação docente, respeitando as condições reais da turma.
